% =====================================================
% part4_eigen.tex ― 第4部:固有値と対角化
% =====================================================
\part{固有値と対角化}

% =====================================================
\chapter{固有値と固有ベクトル}\label{chap:eigen}
% =====================================================

\begin{goal}
\begin{itemize}
\item 固有値・固有ベクトルの定義を、「変換が向きを変えない特別な方向」
という幾何的意味とともに理解する。
\item 固有多項式 $\det(\lambda E - A) = 0$ を解いて固有値を、
同次方程式を掃き出して固有ベクトル(固有空間)を、確実に計算できるようになる。
\item 固有値と正則性・三角行列・トレース・行列式の関係を使いこなし、
「固有値は座標によらない変換そのものの属性」であることを説明できるようになる。
\end{itemize}
\end{goal}

\section{動機 ― 変換が「向きを変えない」方向}

第III部の最後の問 \ref{q:repmat} で、不思議な実験をした。
変換 $A = \mat{3 & 1 \\ 1 & 3}$ を、基底
$\vect{e}'_1 = \mat{1 \\ 1}$, $\vect{e}'_2 = \mat{1 \\ -1}$ で見直すと、
表現行列が対角行列 $\mat{4 & 0 \\ 0 & 2}$ になった。
種を明かそう。実はこの $2$ 本のベクトルは
\[
A\mat{1 \\ 1} = \mat{4 \\ 4} = 4\mat{1 \\ 1},
\qquad
A\mat{1 \\ -1} = \mat{2 \\ -2} = 2\mat{1 \\ -1}
\]
を満たす。つまり、\textbf{変換 $A$ で向きを変えず、ただ定数倍されるだけの
特別な方向}だったのである。一般のベクトルは $A$ で向きも長さも変わるのに、
この $2$ 方向だけは「軸」のように振る舞う。
変換に内在するこの特別な方向こそ、線形代数最重要の概念である。

\section{定義}

\begin{teigi}[固有値・固有ベクトル]
$n$ 次正方行列 $A$ に対して、
\[
A\vect{v} = \lambda \vect{v}, \qquad \vect{v} \neq \vect{0}
\]
を満たす数 $\lambda$ とベクトル $\vect{v}$ が存在するとき、
$\lambda$ を $A$ の\textbf{固有値}、$\vect{v}$ を
固有値 $\lambda$ に属する\textbf{固有ベクトル}という。
\end{teigi}

\begin{chui}[$\vect{v} \neq \vect{0}$ と $\lambda = 0$]
$\vect{v} \neq \vect{0}$ の条件は本質的である
($\vect{v} = \vect{0}$ ならどんな $\lambda$ でも式が成り立ってしまい、
何の情報もない)。一方、\textbf{固有値} $\lambda = 0$ は許される:
$A\vect{v} = \vect{0}$ となる $\vect{v} \neq \vect{0}$ の存在を意味し、
後述の定理 \ref{thm:eigen-zero} のとおり「$A$ が正則でない」ことと同値になる。
「固有\textbf{ベクトル}は零を除外、固有\textbf{値}は零でもよい」と
区別して覚えること。
\end{chui}

\section{固有値の求め方 ― 固有方程式}

定義の式を移項すると
\[
(\lambda E - A)\vect{v} = \vect{0}, \qquad \vect{v} \neq \vect{0}
\]
となる。これは「同次方程式が非自明な解をもつ」条件である。

\begin{teiri}[固有値の特徴づけ]\label{thm:eigen-char}
$\lambda$ が $A$ の固有値であることと、
\[
\det(\lambda E - A) = 0
\]
は同値である。
\end{teiri}

\noindent\textbf{(証明)}\quad
$\lambda$ が固有値
$\iff$ $(\lambda E - A)\vect{v} = \vect{0}$ が非自明な解 $\vect{v} \neq \vect{0}$
をもつ
$\iff$ $\lambda E - A$ が正則でない
(正則なら解は $\vect{v} = \vect{0}$ のみ、正則でなければ自由度が $1$ 以上)
$\iff$ $\det(\lambda E - A) = 0$
(第II部の定理 \ref{thm:det-nonzero})。 $\blacksquare$

左辺 $\det(\lambda E - A)$ は $\lambda$ の $n$ 次多項式になる。

\begin{teigi}[固有多項式]
$\varphi_A(\lambda) = \det(\lambda E - A)$ を $A$ の\textbf{固有多項式}、
方程式 $\varphi_A(\lambda) = 0$ を\textbf{固有方程式}という。
$A$ の固有値とは、固有方程式の解にほかならない。
\end{teigi}

求める手順は完全に機械化される:
\begin{enumerate}
\item $\det(\lambda E - A) = 0$ を解いて固有値 $\lambda$ を求める。
\item 各 $\lambda$ について、同次方程式 $(\lambda E - A)\vect{v} = \vect{0}$ を
掃き出し法で解き、固有ベクトルを求める。
\end{enumerate}
固有値は行列式が、固有ベクトルは掃き出し法が担当する
――第I・II部の道具立てが、ここで一斉に稼働する。

\begin{rei}
$A = \mat{3 & 1 \\ 1 & 3}$ の固有値・固有ベクトルを求める。
\[
\varphi_A(\lambda) = \det \mat{\lambda - 3 & -1 \\ -1 & \lambda - 3}
= (\lambda - 3)^2 - 1 = (\lambda - 4)(\lambda - 2)
\]
より固有値は $\lambda = 4, 2$。

$\lambda = 4$ のとき $(4E - A)\vect{v} = \vect{0}$ は
$\mat{1 & -1 \\ -1 & 1}\vect{v} = \vect{0}$ となり、
解は $\vect{v} = t\mat{1 \\ 1}$($t \neq 0$)。

$\lambda = 2$ のとき同様に $\vect{v} = t\mat{1 \\ -1}$($t \neq 0$)。

冒頭の「特別な方向」が、計算で再発見された。
\end{rei}

\begin{chui}[検算のすすめ]
固有ベクトルを求めたら、必ず $A\vect{v} = \lambda\vect{v}$ に戻して
検算すること。上の例なら
$A\mat{1 \\ 1} = \mat{4 \\ 4} = 4\mat{1 \\ 1}$ と、一行で確かめられる。
また $(\lambda E - A)\vect{v} = \vect{0}$ を解いたとき
\textbf{解が $\vect{v} = \vect{0}$ しか出なかったら、
固有値の計算が間違っている}(固有値なら非自明解が必ずあるはず)。
掃き出しの途中で行が全部消えることが、むしろ正しさの証拠である。
\end{chui}

\section{固有値がすぐ分かる行列}

\begin{teiri}[$0$ が固有値 $\iff$ 正則でない]\label{thm:eigen-zero}
$0$ が $A$ の固有値であることと、$A$ が正則でないことは同値である。
\end{teiri}

\noindent\textbf{(証明)}\quad
定理 \ref{thm:eigen-char} で $\lambda = 0$ とすると、
$0$ が固有値 $\iff \det(0 \cdot E - A) = \det(-A) = (-1)^n \det A = 0
\iff \det A = 0 \iff A$ は正則でない
(定理 \ref{thm:det-nonzero})。 $\blacksquare$

\begin{teiri}[三角行列の固有値]\label{thm:eigen-tri}
上三角行列(下三角行列・対角行列も同様)の固有値は、
対角成分そのものである。
\end{teiri}

\noindent\textbf{(証明)}\quad
$A$ が上三角なら $\lambda E - A$ も上三角で、その対角成分は
$\lambda - a_{11}, \dots, \lambda - a_{nn}$ である。三角行列の行列式は
対角成分の積だから(定理 \ref{thm:det-tri})
\[
\varphi_A(\lambda) = (\lambda - a_{11})(\lambda - a_{22})\cdots(\lambda - a_{nn}).
\]
この根は $a_{11}, \dots, a_{nn}$ である。 $\blacksquare$

ただし、\textbf{固有ベクトルまで自明とは限らない}
(標準基底とは限らない。例題 \ref{rei:eigen-tri} で確かめる)。

\begin{teiri}[トレース・行列式と固有値($2$ 次)]\label{thm:eigen-trdet2}
$2$ 次正方行列 $A = \mat{a & b \\ c & d}$ の固有多項式は
\[
\varphi_A(\lambda) = \lambda^2 - (\mathrm{tr}\,A)\,\lambda + \det A
\]
であり、固有値 $\lambda_1, \lambda_2$(重解も込めて)について
\[
\lambda_1 + \lambda_2 = \mathrm{tr}\,A, \qquad
\lambda_1 \lambda_2 = \det A
\]
が成り立つ。
\end{teiri}

\noindent\textbf{(証明)}\quad
\[
\varphi_A(\lambda) = \det\mat{\lambda - a & -b \\ -c & \lambda - d}
= (\lambda - a)(\lambda - d) - bc
= \lambda^2 - (a + d)\lambda + (ad - bc)
\]
で前半が従う。後半は $2$ 次方程式の解と係数の関係
($\varphi_A(\lambda) = (\lambda - \lambda_1)(\lambda - \lambda_2)$ を
展開して係数を比べる)である。 $\blacksquare$

固有値を求めたら「和がトレース、積が行列式」で検算する習慣を
つけるとよい。同じ関係は $n$ 次でも成り立つ(深掘り参照)。

\section{固有空間}

固有値 $\lambda$ に属する固有ベクトル全体に $\vect{0}$ を付け加えた集合
\[
W_\lambda = \{\, \vect{v} \mid A\vect{v} = \lambda\vect{v} \,\}
= \mathrm{Ker}(\lambda E - A)
\]
は、同次方程式の解空間だから\textbf{部分空間}になる
(第 \ref{chap:linmap} 章の言葉でいえば、線形写像
$\vect{x} \mapsto (\lambda E - A)\vect{x}$ の核である)。
これを固有値 $\lambda$ の\textbf{固有空間}という。
固有ベクトルは $1$ 本ではなく、空間ごと存在するのである。
固有空間の次元 $\dim W_\lambda$ は、自由度定理により
$n - \mathrm{rank}(\lambda E - A)$ で計算できる。

\begin{itembox}[l]{深掘り:固有値は「変換そのもの」の属性である}
第III部の深掘りで、相似な行列($A' = P^{-1}AP$)は同じ変換を
別の座標系から見たものだと述べた。固有多項式は相似で変わらない:
\[
\det(\lambda E - P^{-1}AP) = \det\bigl(P^{-1}(\lambda E - A)P\bigr)
= \det P^{-1} \det(\lambda E - A) \det P = \varphi_A(\lambda).
\]
よって\textbf{固有値は相似不変}、すなわち座標の取り方によらない
「変換そのものの属性」である。さらに、固有多項式を
$\varphi_A(\lambda) = (\lambda - \lambda_1)\cdots(\lambda - \lambda_n)$ と
因数分解して係数を比較すると(定理 \ref{thm:eigen-trdet2} は $n = 2$ の場合)、
\[
\lambda_1 + \cdots + \lambda_n = \mathrm{tr}\,A,
\qquad
\lambda_1 \cdots \lambda_n = \det A
\]
が分かる。第III部で「相似不変量」として登場したトレースと行列式の正体は、
\textbf{固有値の和と積}だったのである。
行列式 = 体積拡大率 = 固有値の積、という連鎖は直観とも整合する:
各固有方向に $\lambda_i$ 倍ずつ伸ばせば、体積は $\lambda_1 \cdots \lambda_n$ 倍になる。
\end{itembox}

\begin{mondai}\label{q:eigen}
$A = \mat{1 & 4 \\ 2 & 3}$ の固有値と、各固有値に属する固有ベクトルを求めよ。
また、固有値の和と積がそれぞれ $\mathrm{tr}\,A$、$\det A$ に
一致することを確かめよ。
\end{mondai}

\section{反例 ― ありがちな誤解を正す}

\begin{hanrei}[固有ベクトルの和は固有ベクトルとは限らない]\label{hanrei:eigen-sum}
「固有ベクトルどうしを足しても固有ベクトル」ではない。
$A = \mat{1 & 0 \\ 0 & 2}$ の固有ベクトル
$\vect{e}_1 = \mat{1 \\ 0}$(固有値 $1$)と
$\vect{e}_2 = \mat{0 \\ 1}$(固有値 $2$)の和
$\vect{v} = \mat{1 \\ 1}$ は、
\[
A\vect{v} = \mat{1 \\ 2}
\]
となり、$\vect{v}$ のどんな定数倍でもない。よって固有ベクトルでない。
和が固有ベクトルに留まるのは、\textbf{同じ固有値}に属する
固有ベクトルどうし(=同じ固有空間の中)だけである。
「固有ベクトル全体」は部分空間にならず、固有空間ごとに分かれて
部分空間をなす、と整理しておこう。
\end{hanrei}

\begin{hanrei}[固有値は足し算・掛け算と相性が悪い]\label{hanrei:eigen-add}
「$A + B$ の固有値は $A$ の固有値と $B$ の固有値の和」は誤りである。
\[
A = \mat{0 & 1 \\ 0 & 0}, \qquad B = \mat{0 & 0 \\ 1 & 0}
\]
はどちらも固有多項式 $\lambda^2$ をもち、固有値は $0$ のみ。
しかし $A + B = \mat{0 & 1 \\ 1 & 0}$ の固有多項式は
$\lambda^2 - 1$ で、固有値は $\pm 1$ である($0 + 0$ ではない)。
積 $AB = \mat{1 & 0 \\ 0 & 0}$ の固有値も $1, 0$ で、
「固有値の積($0 \cdot 0 = 0$)」からは説明できない。
固有値は行列の和・積とは(一般には)連動しない。
連動するのは\textbf{同じ $A$ のべき}である:$A\vect{v} = \lambda\vect{v}$
ならば $A^2\vect{v} = \lambda^2\vect{v}$(問 \ref{q:eigen-ex8})。
\end{hanrei}

\section{例題}

固有値・固有ベクトルの計算と検証、そして性質の証明を練習する。
以下は\textbf{解答つきの例題}である。手を動かして追ってほしい。

\begin{rei}[例題・基本(固有値と固有ベクトルの計算)]\label{rei:eigen-basic}
$A = \mat{2 & 1 \\ 1 & 2}$ の固有値と固有ベクトルを求めよ。\\
\textbf{解.}\ 固有多項式は
\[
\varphi_A(\lambda) = \det\mat{\lambda - 2 & -1 \\ -1 & \lambda - 2}
= (\lambda - 2)^2 - 1 = (\lambda - 3)(\lambda - 1)
\]
より固有値は $\lambda = 3, 1$
(検算:和 $3 + 1 = 4 = \mathrm{tr}\,A$、積 $3 \cdot 1 = 3 = \det A$)。

$\lambda = 3$:$(3E - A)\vect{v} = \mat{1 & -1 \\ -1 & 1}\vect{v} = \vect{0}$
より $x = y$、$\vect{v} = t\mat{1 \\ 1}$($t \neq 0$)。

$\lambda = 1$:$(E - A)\vect{v} = \mat{-1 & -1 \\ -1 & -1}\vect{v} = \vect{0}$
より $x + y = 0$、$\vect{v} = t\mat{1 \\ -1}$($t \neq 0$)。

例題 \ref{rei:repmat-change-ex} で「この基底なら表現行列が
$\mathrm{diag}(3, 1)$ になる」と観察した、その基底が再登場した。
\end{rei}

\begin{rei}[例題・基本(三角行列)]\label{rei:eigen-tri}
上三角行列 $A = \mat{3 & 1 \\ 0 & 2}$ の固有値と固有ベクトルを求めよ。\\
\textbf{解.}\ 定理 \ref{thm:eigen-tri} より、固有値は対角成分の
$\lambda = 3, 2$ である(固有方程式
$(\lambda - 3)(\lambda - 2) = 0$ を解いても同じ)。

$\lambda = 3$:$(3E - A)\vect{v} = \mat{0 & -1 \\ 0 & 1}\vect{v} = \vect{0}$
より $y = 0$、$\vect{v} = t\mat{1 \\ 0}$。

$\lambda = 2$:$(2E - A)\vect{v} = \mat{-1 & -1 \\ 0 & 0}\vect{v} = \vect{0}$
より $x + y = 0$、$\vect{v} = t\mat{1 \\ -1}$。

固有値は対角成分から即座に読めたが、固有ベクトルの一方
$\mat{1 \\ -1}$ は標準基底ベクトルではない。
\textbf{「三角行列だから固有ベクトルも自明」ではない}点に注意。
\end{rei}

\begin{rei}[例題・基本(固有ベクトルかどうかの判定)]\label{rei:eigen-check}
$A = \mat{1 & 2 \\ 3 & 2}$ について、
$\vect{v}_1 = \mat{1 \\ 1}$, $\vect{v}_2 = \mat{2 \\ 3}$ は
$A$ の固有ベクトルか。固有ベクトルであれば固有値を答えよ。\\
\textbf{解.}\ 定義に従い、$A$ を掛けて定数倍になるかを見る。
\[
A\vect{v}_1 = \mat{1 + 2 \\ 3 + 2} = \mat{3 \\ 5}
\]
は $\mat{1 \\ 1}$ の定数倍ではない($3 \neq 5$)ので、
$\vect{v}_1$ は固有ベクトルで\textbf{ない}。
\[
A\vect{v}_2 = \mat{2 + 6 \\ 6 + 6} = \mat{8 \\ 12} = 4\mat{2 \\ 3}
\]
より、$\vect{v}_2$ は固有値 $4$ に属する固有ベクトルで\textbf{ある}。
判定に固有方程式は不要である――掛けてみればよい。
\end{rei}

\begin{rei}[例題・標準($3$ 次行列)]\label{rei:eigen-3x3}
$A = \mat{1 & 0 & 2 \\ 0 & 3 & 0 \\ 2 & 0 & 1}$ の固有値と、
各固有空間の基底を求めよ。\\
\textbf{解.}\ 固有多項式を第 $2$ 行で余因子展開すると
\[
\varphi_A(\lambda) = \det\mat{\lambda - 1 & 0 & -2 \\ 0 & \lambda - 3 & 0 \\
-2 & 0 & \lambda - 1}
= (\lambda - 3)\bigl\{(\lambda - 1)^2 - 4\bigr\}
= (\lambda - 3)^2(\lambda + 1).
\]
固有値は $\lambda = 3$(重解)と $\lambda = -1$
(検算:和 $3 + 3 - 1 = 5 = \mathrm{tr}\,A$)。

$\lambda = 3$:$(3E - A)\vect{v} = \mat{2 & 0 & -2 \\ 0 & 0 & 0 \\
-2 & 0 & 2}\vect{v} = \vect{0}$ は $x = z$($y$ は自由)を与えるから、
\[
W_3 = \left\{ s\mat{1 \\ 0 \\ 1} + t\mat{0 \\ 1 \\ 0} \right\}
\]
で、基底は $\mat{1 \\ 0 \\ 1}, \mat{0 \\ 1 \\ 0}$($2$ 次元)。

$\lambda = -1$:$(-E - A)\vect{v} = \mat{-2 & 0 & -2 \\ 0 & -4 & 0 \\
-2 & 0 & -2}\vect{v} = \vect{0}$ は $x = -z$, $y = 0$ を与えるから、
基底は $\mat{1 \\ 0 \\ -1}$($1$ 次元)。

重解の固有値 $3$ に対して固有空間が $2$ 次元ある
――この「重複度どうしの一致」が次章の対角化可能性の鍵になる。
\end{rei}

\begin{rei}[例題・標準(固有値 $0$ と正則性)]\label{rei:eigen-zero-ex}
$A = \mat{2 & 6 \\ 1 & 3}$ の固有値と固有ベクトルを求め、
$0$ が固有値であることと $\det A = 0$ の関係を確かめよ。\\
\textbf{解.}\ $\varphi_A(\lambda) = (\lambda - 2)(\lambda - 3) - 6
= \lambda^2 - 5\lambda = \lambda(\lambda - 5)$ より固有値は
$\lambda = 0, 5$。

$\lambda = 0$:$A\vect{v} = \vect{0}$、すなわち $x + 3y = 0$
($2$ 行は比例)より $\vect{v} = t\mat{-3 \\ 1}$。

$\lambda = 5$:$(5E - A)\vect{v} = \mat{3 & -6 \\ -1 & 2}\vect{v} = \vect{0}$
より $x = 2y$、$\vect{v} = t\mat{2 \\ 1}$
(検算:$A\mat{2 \\ 1} = \mat{10 \\ 5} = 5\mat{2 \\ 1}$)。

実際 $\det A = 6 - 6 = 0$ で $A$ は正則でなく、
定理 \ref{thm:eigen-zero} のとおり $0$ が固有値になっている。
固有値 $0$ の固有空間 $W_0 = \mathrm{Ker}\,A$ は、
同次方程式 $A\vect{x} = \vect{0}$ の解空間そのものである。
\end{rei}

\begin{rei}[例題・標準(証明:射影の固有値)]\label{rei:eigen-proj}
$P = \dfrac{1}{2}\mat{1 & 1 \\ 1 & 1}$ は、直線 $y = x$ への正射影
(各点から直線に下ろした垂線の足へうつす変換)を表す。
$P$ の固有値と固有ベクトルを求め、幾何的に解釈せよ。\\
\textbf{解.}\ $\varphi_P(\lambda) = \left(\lambda - \dfrac{1}{2}\right)^2
- \dfrac{1}{4} = \lambda^2 - \lambda = \lambda(\lambda - 1)$ より
固有値は $\lambda = 1, 0$。

$\lambda = 1$:$\vect{v} = t\mat{1 \\ 1}$
(検算:$P\mat{1 \\ 1} = \mat{1 \\ 1}$)。
直線 $y = x$ 上の点は射影で\textbf{動かない}(固有値 $1$)。

$\lambda = 0$:$\vect{v} = t\mat{1 \\ -1}$
(検算:$P\mat{1 \\ -1} = \vect{0}$)。
直線に垂直な方向は射影で\textbf{つぶれる}(固有値 $0$)。

「動かない方向」と「つぶれる方向」という射影の幾何的本性が、
固有値 $1$ と $0$ にそのまま現れている。なお $P^2 = P$
(2度射影しても同じ)であり、固有値も $1^2 = 1$, $0^2 = 0$ と
つじつまが合っている。
\end{rei}

\begin{matome}
\begin{itemize}
\item \textbf{固有ベクトル}=変換で向きを変えない($\lambda$ 倍される
だけの)零でないベクトル。\textbf{固有値} $\lambda$ は $0$ でもよく、
$0$ が固有値 $\iff A$ が正則でない(定理 \ref{thm:eigen-zero})。
\item 固有値は\textbf{固有方程式} $\det(\lambda E - A) = 0$ の解
(定理 \ref{thm:eigen-char})、固有ベクトルは同次方程式
$(\lambda E - A)\vect{v} = \vect{0}$ の非自明解。
求めたら $A\vect{v} = \lambda\vect{v}$ で検算する。
\item 三角行列の固有値は対角成分(定理 \ref{thm:eigen-tri})。
ただし固有ベクトルは別途計算が要る。
\item 固有値の\textbf{和はトレース、積は行列式}
(定理 \ref{thm:eigen-trdet2}、$n$ 次でも成立)。検算に活用できる。
\item 固有値 $\lambda$ の\textbf{固有空間}
$W_\lambda = \mathrm{Ker}(\lambda E - A)$ は部分空間で、
次元は $n - \mathrm{rank}(\lambda E - A)$。
異なる固有値の固有ベクトルどうしの和は固有ベクトルにならない
(反例 \ref{hanrei:eigen-sum})。
\item 固有多項式は相似不変。固有値は座標の取り方によらない
\textbf{変換そのものの属性}である。
\end{itemize}
\end{matome}

\section*{章末問題}

\begin{mondai}[基本]\label{q:eigen-ex1}
$A = \mat{5 & 4 \\ 2 & 3}$ の固有値と、各固有値に属する固有ベクトルを
求めよ。また固有値の和と積を $\mathrm{tr}\,A$, $\det A$ と見比べよ。
\end{mondai}

\begin{mondai}[基本]\label{q:eigen-ex2}
上三角行列 $A = \mat{2 & 3 \\ 0 & 5}$ の固有値を対角成分から読み取り、
各固有値の固有ベクトルを求めよ。
\end{mondai}

\begin{mondai}[基本]\label{q:eigen-ex3}
$A = \mat{1 & 2 \\ 3 & 2}$ について、次のベクトルが $A$ の固有ベクトル
かどうかを判定し、固有ベクトルの場合は固有値を答えよ。
\[
\vect{v}_1 = \mat{1 \\ 1}, \qquad
\vect{v}_2 = \mat{2 \\ 3}, \qquad
\vect{v}_3 = \mat{1 \\ -1}
\]
\end{mondai}

\begin{mondai}[基本]\label{q:eigen-ex4}
対角行列 $A = \mat{2 & 0 \\ 0 & -1}$ の固有値と固有空間を求めよ。
一般に、対角行列 $\mathrm{diag}(a_1, \dots, a_n)$ の固有値と、
標準基底の関係を述べよ。
\end{mondai}

\begin{mondai}[標準]\label{q:eigen-ex5}
$A = \mat{2 & 1 & 1 \\ 0 & 1 & 0 \\ 0 & 1 & 3}$ の固有値と、
各固有空間の基底を求めよ(上三角に近い形なので、固有多項式は
第 $1$ 列での余因子展開が速い)。
\end{mondai}

\begin{mondai}[標準]\label{q:eigen-ex6}
$90^\circ$ 回転行列 $R = \mat{0 & -1 \\ 1 & 0}$ には実数の固有値が
存在しないことを、固有方程式を用いて示せ。
また、その事実の幾何的な理由を一行で述べよ。
\end{mondai}

\begin{mondai}[標準]\label{q:eigen-ex7}
$A = \mat{1 & 2 \\ 2 & 4}$ について、次に答えよ。\\
　(1) $\det A$ を計算し、$0$ が固有値であることを予言せよ。\\
　(2) 固有値と固有ベクトルをすべて求めよ。\\
　(3) 固有値 $0$ の固有空間は、第II部で学んだどの空間と一致するか。
\end{mondai}

\begin{mondai}[標準]\label{q:eigen-ex8}
$A\vect{v} = \lambda\vect{v}$($\vect{v} \neq \vect{0}$)のとき、
次を示せ。\\
　(1) $A^2\vect{v} = \lambda^2\vect{v}$(したがって $\vect{v}$ は
$A^2$ の固有値 $\lambda^2$ の固有ベクトル)。\\
　(2) $A$ が正則ならば $\lambda \neq 0$ であり、
$A^{-1}\vect{v} = \dfrac{1}{\lambda}\vect{v}$。\\
　(3) 例題 \ref{rei:eigen-basic} の $A = \mat{2 & 1 \\ 1 & 2}$ について、
$A^2$ の固有値を計算なしで予言し、$A^2$ を実際に計算して確かめよ。
\end{mondai}

\begin{mondai}[発展]\label{q:eigen-ex9}
$A = \mat{0 & 1 \\ 1 & 1}$ について、次に答えよ。\\
　(1) 固有値 $\varphi = \dfrac{1 + \sqrt{5}}{2}$,
$\psi = \dfrac{1 - \sqrt{5}}{2}$ を固有方程式から求めよ。\\
　(2) 各固有値 $\lambda$ に対し、$\mat{1 \\ \lambda}$ が固有ベクトルに
なることを、$\lambda^2 = \lambda + 1$ を用いて確かめよ。
(この行列は第 \ref{chap:power} 章でフィボナッチ数列の解析に使われる。)
\end{mondai}

\begin{mondai}[発展]\label{q:eigen-ex10}
鏡映(直線 $y = x \tan\theta$ に関する折り返し)を表す行列
$S = \mat{\cos 2\theta & \sin 2\theta \\ \sin 2\theta & -\cos 2\theta}$
について、固有値が $\theta$ によらず $1$ と $-1$ であることを
固有方程式から示せ。また、それぞれの固有ベクトルの方向を
幾何的に説明せよ(計算では $\theta = 30^\circ$ など特定の場合を
確かめればよい)。
\end{mondai}

\begin{mondai}[証明]\label{q:eigen-ex11}
$A$ と転置行列 ${}^t\!A$ の固有多項式が一致する(したがって固有値も
一致する)ことを証明せよ
(ヒント:$\lambda E - {}^t\!A = {}^t(\lambda E - A)$ と、
転置で行列式が変わらないこと(定理 \ref{thm:det-props})を使う)。
\end{mondai}

\begin{mondai}[証明]\label{q:eigen-ex12}
$B = P^{-1}AP$ とする($P$ は正則)。次を証明せよ。\\
　(1) $A$ と $B$ の固有多項式は一致する。\\
　(2) $\vect{v}$ が $A$ の固有値 $\lambda$ の固有ベクトルならば、
$P^{-1}\vect{v}$ は $B$ の固有値 $\lambda$ の固有ベクトルである。
(相似=座標変換で、固有値は不変、固有ベクトルは座標が変わるだけ。)
\end{mondai}

\bigskip
{\small
\noindent\textbf{【章末問題の方針】}\quad
まず自力で取り組み、行き詰まったら下の方針を見よ。記述による完全解答は
巻末「問の解答」にある。
\begin{itemize}
\item \textbf{問 \ref{q:eigen-ex1}}:固有方程式を解き、各固有値で
$(\lambda E - A)\vect{v} = \vect{0}$ を掃き出す。和と積で検算。
\item \textbf{問 \ref{q:eigen-ex2}}:固有値は定理 \ref{thm:eigen-tri}。
固有ベクトルは各固有値で同次方程式を解く。
\item \textbf{問 \ref{q:eigen-ex3}}:$A$ を掛けて、結果がもとの
ベクトルの定数倍かどうかを見るだけでよい。
\item \textbf{問 \ref{q:eigen-ex4}}:$A\vect{e}_1$, $A\vect{e}_2$ を
計算してみる。
\item \textbf{問 \ref{q:eigen-ex5}}:第 $1$ 列で余因子展開すると
固有多項式が $(\lambda - 2)(\lambda - 1)(\lambda - 3)$ と
因数分解された形で得られる。固有ベクトルは各固有値で掃き出す。
\item \textbf{問 \ref{q:eigen-ex6}}:判別式が負になることを見る。
幾何的には「すべての方向が回されてしまう」。
\item \textbf{問 \ref{q:eigen-ex7}}:(3) $W_0 = \mathrm{Ker}\,A$ は
同次方程式の解空間である。
\item \textbf{問 \ref{q:eigen-ex8}}:(1) は $A$ を $2$ 回掛ける。
(2) は両辺に $A^{-1}$ を掛ける($\lambda = 0$ だと $A\vect{v} = \vect{0}$
で正則性に矛盾)。
\item \textbf{問 \ref{q:eigen-ex9}}:(2) は
$A\mat{1 \\ \lambda} = \mat{\lambda \\ 1 + \lambda}$ を計算し、
$1 + \lambda = \lambda^2$ で書き換える。
\item \textbf{問 \ref{q:eigen-ex10}}:$\mathrm{tr}\,S = 0$,
$\det S = -\cos^2 2\theta - \sin^2 2\theta = -1$ から固有方程式は
$\lambda^2 - 1 = 0$。鏡の上の方向は不動、垂直方向は反転。
\item \textbf{問 \ref{q:eigen-ex11}}:単位行列は転置しても変わらない。
\item \textbf{問 \ref{q:eigen-ex12}}:(1) は深掘りの計算をなぞる。
(2) は $B(P^{-1}\vect{v})$ を定義どおり計算する。
$P^{-1}\vect{v} \neq \vect{0}$ の確認も忘れずに。
\end{itemize}
}

% =====================================================
\chapter{対角化}\label{chap:diag}
% =====================================================

\begin{goal}
\begin{itemize}
\item 対角化の定義と基本定理(独立な固有ベクトル $n$ 本 $\iff$ 対角化可能)
を証明つきで理解し、$2$ 次・$3$ 次の行列を実際に対角化できるようになる。
\item 「相異なる固有値の固有ベクトルは線形独立」の証明を理解し、
固有値がすべて異なる行列が必ず対角化できる理由を説明できるようになる。
\item 固有値が重複する場合に、固有空間の次元(幾何的重複度)を調べて
対角化可能かどうかを判定できるようになる。
\end{itemize}
\end{goal}

\section{動機 ― 最良の座標系を探す}

第III部の終わりに立てた問いを思い出そう。座標(基底)の取り方は自由
なのだから、\textbf{表現行列が最も簡単になる基底}を選べばよいのでは
ないか?――前章で、その「最良の基底」の候補が見つかった。固有ベクトル
である。固有ベクトルの方向では、変換はただの $\lambda$ 倍であり、
これ以上簡単な作用はない。

もし空間全体を固有ベクトルの基底で張り尽くせたら、どの基底方向でも
変換は定数倍、すなわち表現行列は\textbf{対角行列}になる。
これが対角化である。本章では、いつそれが可能で、いつ不可能なのかを
完全に見極める。

\section{対角化の定義と基本定理}

\begin{teigi}[対角化]
$n$ 次正方行列 $A$ に対して、正則行列 $P$ をうまく選んで
\[
P^{-1} A P = \mat{\lambda_1 & & \\ & \ddots & \\ & & \lambda_n}
\]
(対角行列)とできるとき、$A$ は\textbf{対角化可能}であるといい、
この操作を\textbf{対角化}という。
\end{teigi}

第III部の言葉でいえば、対角化とは
\textbf{表現行列が対角行列になるような基底を見つけること}である。
そして、その基底の正体はもう分かっている。

\begin{teiri}[対角化の基本定理]\label{thm:diag-main}
$n$ 次正方行列 $A$ が対角化可能であるための必要十分条件は、
\textbf{$A$ の固有ベクトルから線形独立な $n$ 本を選べる}ことである。
このとき、固有ベクトル $\vect{p}_1, \dots, \vect{p}_n$ を列に並べた
$P = (\vect{p}_1 \cdots \vect{p}_n)$ により
$P^{-1}AP$ は対角行列になり、対角成分には対応する固有値が
\textbf{並べた順に}現れる。
\end{teiri}

\noindent\textbf{(証明)}\quad
(十分性)独立な固有ベクトル $\vect{p}_1, \dots, \vect{p}_n$
($A\vect{p}_j = \lambda_j\vect{p}_j$)があるとする。
行列の積の列ごとの計算により
\[
AP = (A\vect{p}_1 \ \cdots \ A\vect{p}_n)
= (\lambda_1 \vect{p}_1 \ \cdots \ \lambda_n \vect{p}_n) = P
\mat{\lambda_1 & & \\ & \ddots & \\ & & \lambda_n} = PD
\]
となる。$P$ は列が線形独立だから正則(第II部の定理 \ref{thm:reg-rank})で、
左から $P^{-1}$ を掛けて $P^{-1}AP = D$。

(必要性)逆に $P^{-1}AP = D$(対角)とすると、両辺に左から $P$ を掛けて
$AP = PD$。第 $j$ 列を比べると $A\vect{p}_j = \lambda_j \vect{p}_j$ であり、
$P$ は正則だから各列 $\vect{p}_j$ は $\vect{0}$ でなく、$n$ 本は線形独立。
すなわち独立な固有ベクトルが $n$ 本とれている。 $\blacksquare$

問 \ref{q:repmat} で起きたことの完全な説明がこれである。
実際の計算手順は次のとおり機械化される。
\begin{enumerate}
\item 固有方程式を解いて固有値をすべて求める。
\item 各固有値の固有空間の基底を掃き出し法で求める。
\item 集めた固有ベクトルが $n$ 本(独立)あれば、列に並べて $P$ を作る。
$P^{-1}AP$ は、\textbf{並べた順に}固有値が並ぶ対角行列 $D$ になる。
\item 検算は $AP = PD$ で行う($P^{-1}$ の計算は不要)。
\end{enumerate}

では、独立な固有ベクトルは $n$ 本そろうのか。
強力な十分条件が一つある。

\section{相異なる固有値なら必ず対角化できる}

\begin{teiri}[相異なる固有値の固有ベクトルは独立]\label{thm:diag-distinct}
$\lambda_1, \dots, \lambda_k$ を $A$ の相異なる固有値、
$\vect{v}_1, \dots, \vect{v}_k$ をそれぞれに属する固有ベクトルとすると、
$\vect{v}_1, \dots, \vect{v}_k$ は線形独立である。
\end{teiri}

\noindent\textbf{(証明)}\quad
本数 $k$ についての数学的帰納法で示す。$k = 1$ のときは、
固有ベクトルは $\vect{0}$ でないから $1$ 本で独立である。

$k - 1$ 本までは正しいとして、$k$ 本の場合を示す。
\[
c_1\vect{v}_1 + c_2\vect{v}_2 + \cdots + c_k\vect{v}_k = \vect{0}
\tag{$*$}
\]
とする。($*$) の両辺に $A$ を掛けると、$A\vect{v}_i = \lambda_i\vect{v}_i$
より
\[
c_1\lambda_1\vect{v}_1 + c_2\lambda_2\vect{v}_2 + \cdots
+ c_k\lambda_k\vect{v}_k = \vect{0}.
\tag{$**$}
\]
一方、($*$) の両辺を $\lambda_k$ 倍すると
\[
c_1\lambda_k\vect{v}_1 + c_2\lambda_k\vect{v}_2 + \cdots
+ c_k\lambda_k\vect{v}_k = \vect{0}.
\]
($**$) からこれを辺々引くと、最後の項が消えて
\[
c_1(\lambda_1 - \lambda_k)\vect{v}_1 + \cdots
+ c_{k-1}(\lambda_{k-1} - \lambda_k)\vect{v}_{k-1} = \vect{0}.
\]
帰納法の仮定により $\vect{v}_1, \dots, \vect{v}_{k-1}$ は独立だから、
係数はすべて $0$:$c_i(\lambda_i - \lambda_k) = 0$($i \le k - 1$)。
固有値は相異なるので $\lambda_i - \lambda_k \neq 0$、よって
$c_1 = \cdots = c_{k-1} = 0$。これを ($*$) に戻すと
$c_k\vect{v}_k = \vect{0}$ で、$\vect{v}_k \neq \vect{0}$ より $c_k = 0$。
以上ですべての係数が $0$ となり、独立が示された。 $\blacksquare$

\begin{teiri}[$n$ 個の相異なる固有値をもてば対角化可能]\label{thm:diag-ndistinct}
$n$ 次正方行列 $A$ が $n$ 個の相異なる実数の固有値をもてば、
$A$ は対角化可能である。
\end{teiri}

\noindent\textbf{(証明)}\quad
各固有値から固有ベクトルを $1$ 本ずつとれば、
定理 \ref{thm:diag-distinct} により独立な $n$ 本になる。
あとは基本定理 \ref{thm:diag-main} による。 $\blacksquare$

固有方程式は $n$ 次方程式だから解は高々 $n$ 個。
「$n$ 個の相異なる固有値」はいわば\textbf{満員御礼}の状況であり、
このとき対角化は自動的に成功する。
問題は、固有方程式が重解をもつ場合である。

\section{固有値が重複する場合}

\begin{teigi}[代数的重複度・幾何的重複度]
固有値 $\lambda$ の、固有方程式の解としての重複度
(固有多項式が $(\lambda' - \lambda)^m$ で割り切れる最大の $m$)を
\textbf{代数的重複度}、固有空間の次元 $\dim W_\lambda$ を
\textbf{幾何的重複度}と呼ぶ。
\end{teigi}

一般に
\[
1 \le (\text{幾何的重複度}) \le (\text{代数的重複度})
\]
が成り立つ(左の不等号は固有ベクトルの存在そのもの。
右の不等号の証明の粗筋は深掘りで述べる)。
対角化可能性は、固有空間の次元だけで完全に判定できる。

\begin{teiri}[対角化可能性の判定]\label{thm:diag-criterion}
$n$ 次正方行列 $A$ の相異なる実数の固有値を
$\lambda_1, \dots, \lambda_k$ とすると、
\[
A \text{ が対角化可能} \iff
\dim W_{\lambda_1} + \cdots + \dim W_{\lambda_k} = n.
\]
固有多項式が実数の範囲で $1$ 次式の積に分解するときは、この条件は
「\textbf{すべての固有値で幾何的重複度 = 代数的重複度}」と言い換えられる。
\end{teiri}

\noindent\textbf{(証明)}\quad
($\Leftarrow$)各固有空間 $W_{\lambda_i}$ の基底を集めて
$n$ 本のベクトルの組を作る。この組が線形独立であることを示す。
線形結合が $\vect{0}$ になったとし、固有値ごとに項をまとめて
$\vect{w}_1 + \cdots + \vect{w}_k = \vect{0}$
($\vect{w}_i \in W_{\lambda_i}$)と書く。もしある $\vect{w}_i$ が
$\vect{0}$ でなければ、$\vect{0}$ でないものだけを残した式は
「相異なる固有値に属する固有ベクトルたちの、係数 $1$ の線形結合が
$\vect{0}$」を意味し、定理 \ref{thm:diag-distinct} に矛盾する。
よってすべての $\vect{w}_i = \vect{0}$ であり、各 $W_{\lambda_i}$ の中では
基底の独立性から係数がすべて $0$ になる。ゆえにこの $n$ 本は独立な
固有ベクトルであり、基本定理 \ref{thm:diag-main} により対角化可能である。

($\Rightarrow$)対角化可能なら、基本定理により独立な固有ベクトル
$n$ 本がとれる。それらを固有値ごとに分ければ、$W_{\lambda_i}$ の中に
$n_i$ 本の独立なベクトルがあり($n_1 + \cdots + n_k = n$)、
$\dim W_{\lambda_i} \ge n_i$。辺々足して
$\sum_i \dim W_{\lambda_i} \ge n$。
逆向きの不等式:各固有空間の基底を集めた
$\sum_i \dim W_{\lambda_i}$ 本の組は、($\Leftarrow$)の議論と同じ理由で
線形独立であり、$\R^n$ の独立系は $n$ 本以下(第 \ref{chap:basis} 章の
交換補題 \ref{thm:exchange})だから
$\sum_i \dim W_{\lambda_i} \le n$。よって等号が成り立つ。 $\blacksquare$

\begin{rei}[対角化できない行列]\label{rei:diag-shear}
$A = \mat{1 & 1 \\ 0 & 1}$ の固有多項式は $(\lambda - 1)^2$ で、
固有値は $1$ のみ(代数的重複度 $2$)。固有空間は
\[
(E - A)\vect{v} = \mat{0 & -1 \\ 0 & 0}\vect{v} = \vect{0}
\]
の解空間で、$\vect{v} = t\mat{1 \\ 0}$ の $1$ 次元(幾何的重複度 $1$)。
次元の和は $1 < 2$ であり、\textbf{対角化できない}。
この行列は「せん断(ずらし)変換」であり、どの方向も
(横軸方向以外は)必ず傾けられてしまう。
対角化できない行列をどこまで簡単にできるか
――その完全な答えが第VI部のジョルダン標準形である。
\end{rei}

\begin{itembox}[l]{深掘り:幾何的重複度 $\le$ 代数的重複度の理由}
固有値 $\lambda_0$ の固有空間 $W_{\lambda_0}$ の次元を $m$ とし、
その基底 $\vect{p}_1, \dots, \vect{p}_m$ を $\R^n$ の基底
$\vect{p}_1, \dots, \vect{p}_m, \dots, \vect{p}_n$ に拡張する
(第 \ref{chap:basis} 章の拡張定理)。この基底での $A$ の表現行列
$A' = P^{-1}AP$ を考えると、最初の $m$ 本は固有ベクトルだから、
$A'$ は左上 $m$ 列が「対角に $\lambda_0$、その下は $0$」という形になる:
\[
A' = \mat{\lambda_0 E_m & * \\ O & B}
\]
($E_m$ は $m$ 次単位行列、$B$ は $(n-m)$ 次正方行列)。
この形の行列の行列式は、第 $1$ 列での余因子展開を $m$ 回繰り返すと
分かるように「左上ブロックの行列式 × 右下ブロックの行列式」になるから、
\[
\varphi_A(\lambda) = \varphi_{A'}(\lambda)
= (\lambda - \lambda_0)^m \det(\lambda E_{n-m} - B).
\]
(固有多項式が相似で不変であることは第 \ref{chap:eigen} 章の深掘りで
示した。)よって $\varphi_A$ は $(\lambda - \lambda_0)^m$ で割り切れ、
代数的重複度は $m$ 以上である。
\textbf{固有空間が大きいほど、固有多項式にもその痕跡が残る}
――というのがこの不等式の中身である。
\end{itembox}

\begin{itembox}[l]{深掘り:回転には実数の固有ベクトルがない ― 複素数の招待}
回転行列 $R_\theta = \mat{\cos\theta & -\sin\theta \\ \sin\theta & \cos\theta}$
($0 < \theta < \pi$)は、平面のすべての方向を回してしまうから、
「向きを変えない実ベクトル」は存在しないはずである。実際、固有方程式
\[
\lambda^2 - 2\cos\theta\,\lambda + 1 = 0
\]
の判別式は $4(\cos^2\theta - 1) < 0$ で、実数解をもたない。
しかし\textbf{複素数まで世界を広げれば}、解は
\[
\lambda = \cos\theta \pm i\sin\theta
\]
と求まる。これはオイラーの公式でいう $e^{\pm i\theta}$ であり、
「複素数の世界では、回転は($e^{i\theta}$ 倍という)伸縮にすぎない」
という深い事実を示している。

スカラーを複素数に取り替えても、本書の理論はそのまま成立する
(複素ベクトル空間)。そして複素数の世界では、代数学の基本定理により
\textbf{固有方程式は重複を込めて必ず $n$ 個の解をもつ}。
「固有値が存在しないかも」という心配が消えるため、
進んだ理論(第VI部以降)は複素数上で展開するのが標準である。
実行列の複素固有値は必ず共役のペア $\lambda, \overline{\lambda}$ で現れる。
\end{itembox}

\begin{mondai}\label{q:diag}
　(1) $A = \mat{0 & 1 \\ 1 & 0}$ を対角化せよ
(固有値・固有ベクトルを求め、$P$ と $P^{-1}AP$ を明示すること)。
この変換は幾何的には何か。\\
　(2) $B = \mat{3 & 1 \\ 0 & 3}$ が対角化可能でないことを、
重複度の言葉で説明せよ。
\end{mondai}

\section{反例 ― ありがちな誤解を正す}

\begin{hanrei}[固有値が重複しても対角化できることはある]\label{hanrei:diag-repeat}
「固有値が重複したら対角化不可能」は誤りである。
最も簡単な反例は単位行列 $E$:固有値は $1$(代数的重複度 $n$)だが、
固有空間は $W_1 = \R^n$ 全体($n$ 次元)で、$E$ ははじめから対角である。
$3$ 次の例では、例題 \ref{rei:eigen-3x3} の行列が固有値 $3$ を重複して
もつのに対角化できる(例題 \ref{rei:diag-3x3})。
定理 \ref{thm:diag-ndistinct} は\textbf{十分条件}であって必要条件ではない。
重複があるときは、固有空間の次元を実際に数えて
判定する(定理 \ref{thm:diag-criterion})。
\end{hanrei}

\begin{hanrei}[$P$ の列の順序と $D$ の順序はセットである]\label{hanrei:diag-order}
$A = \mat{3 & 1 \\ 1 & 3}$ は、固有値 $4, 2$ と固有ベクトル
$\mat{1 \\ 1}, \mat{1 \\ -1}$ をもつ(第 \ref{chap:eigen} 章)。ここで
\[
P = \mat{1 & 1 \\ -1 & 1} \quad
\Bigl(\text{列の順序:} \mat{1 \\ -1},\ \mat{1 \\ 1}\Bigr)
\]
と並べたのに $D = \mat{4 & 0 \\ 0 & 2}$ とするのは\textbf{誤り}である。
第 $1$ 列 $\mat{1 \\ -1}$ は固有値 $2$ の固有ベクトルだから、
正しくは $D = \mat{2 & 0 \\ 0 & 4}$。実際
$AP = \mat{2 & 4 \\ -2 & 4}$ と
$P\mat{2 & 0 \\ 0 & 4} = \mat{2 & 4 \\ -2 & 4}$ は一致するが、
$P\mat{4 & 0 \\ 0 & 2}$ は $\mat{4 & 2 \\ -4 & 2}$ で一致しない。
\textbf{$D$ の第 $j$ 対角成分は、$P$ の第 $j$ 列の固有値}
――順序の対応を崩してはならない。
\end{hanrei}

\section{例題}

対角化の実行・判定・応用を練習する。
以下は\textbf{解答つきの例題}である。手を動かして追ってほしい。

\begin{rei}[例題・基本(対角化の実行)]\label{rei:diag-basic}
$A = \mat{1 & 2 \\ 2 & 1}$ を対角化せよ。\\
\textbf{解.}\ $\varphi_A(\lambda) = (\lambda - 1)^2 - 4
= (\lambda - 3)(\lambda + 1)$ より固有値は $3, -1$
(相異なるので対角化可能と即断できる。定理 \ref{thm:diag-ndistinct})。

$\lambda = 3$:$\mat{2 & -2 \\ -2 & 2}\vect{v} = \vect{0}$ より
$\vect{v} = t\mat{1 \\ 1}$。
$\lambda = -1$:$\mat{-2 & -2 \\ -2 & -2}\vect{v} = \vect{0}$ より
$\vect{v} = t\mat{1 \\ -1}$。

$P = \mat{1 & 1 \\ 1 & -1}$ と並べれば
\[
P^{-1}AP = \mat{3 & 0 \\ 0 & -1}.
\]
検算($AP = PD$):
$AP = \mat{3 & -1 \\ 3 & 1}$、
$PD = \mat{1 & 1 \\ 1 & -1}\mat{3 & 0 \\ 0 & -1} = \mat{3 & -1 \\ 3 & 1}$
で一致。
\end{rei}

\begin{rei}[例題・基本(対角化可能性の判定)]\label{rei:diag-judge}
次の行列は対角化可能か。
\[
(1)\ \mat{1 & 1 \\ 0 & 2} \qquad (2)\ \mat{2 & 1 \\ 0 & 2}
\]
\textbf{解.}\ (1) 三角行列だから固有値は対角成分の $1, 2$
(定理 \ref{thm:eigen-tri})。$2$ 次で相異なる $2$ つの固有値をもつから、
\textbf{固有ベクトルを求めるまでもなく}対角化可能である
(定理 \ref{thm:diag-ndistinct})。

(2) 固有値は $2$ のみ(代数的重複度 $2$)。固有空間は
$\mat{0 & -1 \\ 0 & 0}\vect{v} = \vect{0}$ の解空間で
$\vect{v} = t\mat{1 \\ 0}$ の $1$ 次元。次元の和が $1 < 2$ なので
対角化不可能である(定理 \ref{thm:diag-criterion})。
同じ三角行列でも、対角成分が重複するときは固有空間の次元を
数えなければ判定できない。
\end{rei}

\begin{rei}[例題・基本($AP = PD$ による検算)]\label{rei:diag-verify}
「$A = \mat{3 & 1 \\ 1 & 3}$ は $P = \mat{1 & 1 \\ 1 & -1}$ により
$D = \mat{4 & 0 \\ 0 & 2}$ に対角化される」という主張
(問 \ref{q:repmat} の結果)を、$P^{-1}$ を計算せずに検算せよ。\\
\textbf{解.}\ $P^{-1}AP = D$ は(両辺に左から $P$ を掛ければ)
$AP = PD$ と同値である。
\[
AP = \mat{3 & 1 \\ 1 & 3}\mat{1 & 1 \\ 1 & -1} = \mat{4 & 2 \\ 4 & -2},
\qquad
PD = \mat{1 & 1 \\ 1 & -1}\mat{4 & 0 \\ 0 & 2} = \mat{4 & 2 \\ 4 & -2}.
\]
一致するから主張は正しい($P$ が正則であること:
$\det P = -2 \neq 0$ も確認)。
逆行列の計算を避けられるので、検算はつねに $AP = PD$ で行うのがよい。
\end{rei}

\begin{rei}[例題・標準($3$ 次行列の対角化)]\label{rei:diag-3x3}
$A = \mat{1 & 0 & 2 \\ 0 & 3 & 0 \\ 2 & 0 & 1}$
(例題 \ref{rei:eigen-3x3})を対角化せよ。\\
\textbf{解.}\ 例題 \ref{rei:eigen-3x3} で求めたとおり、固有値は
$3$(代数的重複度 $2$)と $-1$ で、固有空間は
\[
W_3 = \left\langle \mat{1 \\ 0 \\ 1}, \mat{0 \\ 1 \\ 0} \right\rangle
\ (2 \text{ 次元}), \qquad
W_{-1} = \left\langle \mat{1 \\ 0 \\ -1} \right\rangle \ (1 \text{ 次元}).
\]
次元の和は $2 + 1 = 3$ だから対角化可能である
(定理 \ref{thm:diag-criterion})。基底を列に並べて
\[
P = \mat{1 & 0 & 1 \\ 0 & 1 & 0 \\ 1 & 0 & -1}, \qquad
P^{-1}AP = \mat{3 & 0 & 0 \\ 0 & 3 & 0 \\ 0 & 0 & -1}.
\]
検算($AP = PD$):$AP$ の各列は
$A\mat{1 \\ 0 \\ 1} = \mat{3 \\ 0 \\ 3}$,
$A\mat{0 \\ 1 \\ 0} = \mat{0 \\ 3 \\ 0}$,
$A\mat{1 \\ 0 \\ -1} = \mat{-1 \\ 0 \\ 1}$ で、
$PD$ の各列($3$ 倍、$3$ 倍、$-1$ 倍)と一致する。
\end{rei}

\begin{rei}[例題・標準(パラメータつきの判定)]\label{rei:diag-param}
$A = \mat{1 & 1 \\ 0 & a}$ が対角化可能であるための
$a$ の条件を求めよ。\\
\textbf{解.}\ 三角行列だから固有値は $1$ と $a$ である。

$a \neq 1$ のとき:相異なる $2$ つの固有値をもつから対角化可能
(定理 \ref{thm:diag-ndistinct})。

$a = 1$ のとき:$A = \mat{1 & 1 \\ 0 & 1}$ はせん断行列
(例 \ref{rei:diag-shear})で、固有空間が $1$ 次元しかなく対角化不可能。

よって求める条件は $a \neq 1$ である。
\textbf{パラメータの特別な値でだけ対角化が壊れる}――この現象は、
固有値の衝突(重解化)が原因である。
\end{rei}

\begin{rei}[例題・標準(証明:べきの対角化)]\label{rei:diag-power}
$A$ が $P$ により $P^{-1}AP = D$(対角)と対角化されるならば、
$A^2$ も同じ $P$ により対角化され、$P^{-1}A^2P = D^2$ となることを示せ。\\
\textbf{解.}\ $P^{-1}AP = D$ の両辺を $2$ 乗すると、
\[
D^2 = (P^{-1}AP)(P^{-1}AP) = P^{-1}A\,(PP^{-1})\,AP = P^{-1}A^2P.
\]
真ん中の $PP^{-1} = E$ が消えるのが要点である。$D^2$ は対角行列
(対角成分は $\lambda_i^2$)だから、$A^2$ は $P$ で対角化される。
同様に繰り返せば $P^{-1}A^nP = D^n$、すなわち
\[
A^n = P D^n P^{-1}
\]
が任意の $n$ で成り立つ。この式が次章の主役になる。
\end{rei}

\begin{matome}
\begin{itemize}
\item \textbf{対角化の基本定理}:$A$ が対角化可能 $\iff$ 独立な
固有ベクトルが $n$ 本とれる。$P$ は固有ベクトルを列に並べた行列で、
$D$ には\textbf{並べた順に}固有値が並ぶ(反例 \ref{hanrei:diag-order})。
検算は $AP = PD$ で。
\item \textbf{相異なる固有値の固有ベクトルは線形独立}
(定理 \ref{thm:diag-distinct}、帰納法+固有値の差で消す証明)。
したがって $n$ 個の相異なる固有値をもてば対角化可能。
\item 固有値が重複する場合は、\textbf{固有空間の次元の和が $n$ に
なるか}で判定する(定理 \ref{thm:diag-criterion})。
言い換えれば、(固有多項式が $1$ 次式の積に分解するとき)
すべての固有値で幾何的重複度=代数的重複度。
\item 重複していても対角化できる例(単位行列など)はある
(反例 \ref{hanrei:diag-repeat})。逆に、せん断
$\mat{1 & 1 \\ 0 & 1}$ は固有空間が足りず対角化できない。
その先はジョルダン標準形(第VI部)の出番である。
\item 実数の固有値が存在しない行列(回転)もある。複素数まで広げれば
固有値は必ず $n$ 個そろう(深掘り)。
\end{itemize}
\end{matome}

\section*{章末問題}

\begin{mondai}[基本]\label{q:diag-ex1}
$A = \mat{5 & 4 \\ 2 & 3}$(問 \ref{q:eigen-ex1} の行列)を対角化せよ
($P$ と $P^{-1}AP$ を明示し、$AP = PD$ で検算すること)。
\end{mondai}

\begin{mondai}[基本]\label{q:diag-ex2}
$A = \mat{2 & 1 \\ 1 & 2}$(例題 \ref{rei:eigen-basic} の行列)を
対角化せよ。
\end{mondai}

\begin{mondai}[基本]\label{q:diag-ex3}
次の行列はそれぞれ対角化可能か。理由とともに判定せよ
(対角化の実行までは不要)。
\[
(1)\ \mat{1 & 1 \\ 0 & 1} \qquad
(2)\ \mat{1 & 1 \\ 0 & 2} \qquad
(3)\ \mat{4 & 0 \\ 0 & 4}
\]
\end{mondai}

\begin{mondai}[基本]\label{q:diag-ex4}
$A = \mat{2 & 1 \\ 1 & 2}$ の固有ベクトルを
$\mat{1 \\ -1}$, $\mat{1 \\ 1}$ の\textbf{順に}列に並べた
$P = \mat{1 & 1 \\ -1 & 1}$ を使うとき、$P^{-1}AP$ はどんな対角行列に
なるか。$AP = PD$ の検算で確かめよ(問 \ref{q:diag-ex2} と
対角成分の順序を見比べること)。
\end{mondai}

\begin{mondai}[標準]\label{q:diag-ex5}
$A = \mat{2 & 1 & 1 \\ 0 & 1 & 0 \\ 0 & 1 & 3}$
(問 \ref{q:eigen-ex5} の行列)を対角化せよ。
\end{mondai}

\begin{mondai}[標準]\label{q:diag-ex6}
$A = \mat{3 & 0 & 0 \\ 0 & 2 & 1 \\ 0 & 0 & 2}$ が対角化可能でないことを、
固有空間の次元を計算して示せ。
\end{mondai}

\begin{mondai}[標準]\label{q:diag-ex7}
固有値 $2$ と $5$ をもち、対応する固有ベクトルがそれぞれ
$\mat{1 \\ 1}$, $\mat{1 \\ -1}$ である $2$ 次正方行列 $A$ を求めよ
($A = PDP^{-1}$ を計算する)。
求めた $A$ について $A\mat{1 \\ 1}$ を計算して検算せよ。
\end{mondai}

\begin{mondai}[標準]\label{q:diag-ex8}
$A = \mat{a & 1 \\ 0 & 2}$ が対角化可能であるための $a$ の条件を求めよ。
\end{mondai}

\begin{mondai}[発展]\label{q:diag-ex9}
正射影行列 $P_0 = \dfrac{1}{2}\mat{1 & 1 \\ 1 & 1}$
(例題 \ref{rei:eigen-proj})を対角化し、
$Q^{-1}P_0Q = \mat{1 & 0 \\ 0 & 0}$ となる $Q$ を求めよ。
さらに、この対角化を用いて $P_0^2 = P_0$(2度射影しても同じ)を
計算で示せ。
\end{mondai}

\begin{mondai}[発展]\label{q:diag-ex10}
$A$ は対角化可能で、ある正の整数 $k$ について $A^k = O$ を満たすとする。
このとき $A = O$ であることを示せ
(対角化できるのに「べきで消える」行列は零行列しかない。
逆にいえば、$A \neq O$ なのに $A^k = O$ となる行列
――第VI部の主役「べき零行列」――は決して対角化できない)。
\end{mondai}

\begin{mondai}[証明]\label{q:diag-ex11}
$A$ が対角化可能ならば、転置行列 ${}^t\!A$ も対角化可能であることを
証明せよ(ヒント:$A = PDP^{-1}$ の両辺を転置し、
${}^t(P^{-1}) = ({}^tP)^{-1}$ を使う)。
\end{mondai}

\begin{mondai}[証明]\label{q:diag-ex12}
相異なる固有値 $\lambda_1 \neq \lambda_2$ に属する固有ベクトル
$\vect{v}_1, \vect{v}_2$ が線形独立であることを、帰納法を使わずに
直接証明せよ(ヒント:$c_1\vect{v}_1 + c_2\vect{v}_2 = \vect{0}$ の
両辺に $A$ を掛けた式と、$\lambda_2$ 倍した式を見比べる。
定理 \ref{thm:diag-distinct} の証明の「最初の一歩」を自分の手で
再現する問題である)。
\end{mondai}

\bigskip
{\small
\noindent\textbf{【章末問題の方針】}\quad
まず自力で取り組み、行き詰まったら下の方針を見よ。記述による完全解答は
巻末「問の解答」にある。
\begin{itemize}
\item \textbf{問 \ref{q:diag-ex1}}:固有値・固有ベクトルは
問 \ref{q:eigen-ex1} で求めてある。並べて $P$ を作り、$AP = PD$ で検算。
\item \textbf{問 \ref{q:diag-ex2}}:例題 \ref{rei:eigen-basic} の結果を
使う。
\item \textbf{問 \ref{q:diag-ex3}}:固有値がすべて異なれば即対角化可能。
重複するときだけ固有空間の次元を数える。(3) はそもそも対角行列。
\item \textbf{問 \ref{q:diag-ex4}}:$D$ の対角には「$P$ に並べた順」で
固有値が入る(反例 \ref{hanrei:diag-order})。
\item \textbf{問 \ref{q:diag-ex5}}:$3$ つの固有値が相異なるので
対角化可能。固有ベクトルは問 \ref{q:eigen-ex5} のもの。
\item \textbf{問 \ref{q:diag-ex6}}:固有値 $2$ の固有空間を掃き出して
次元を数える。$1 + 1 < 3$ になるはず。
\item \textbf{問 \ref{q:diag-ex7}}:$A = PDP^{-1}$。$2 \times 2$ の
逆行列公式を使う。
\item \textbf{問 \ref{q:diag-ex8}}:三角行列なので固有値は $a$ と $2$。
$a = 2$ の場合だけ固有空間を調べる。
\item \textbf{問 \ref{q:diag-ex9}}:固有値 $1, 0$ の固有ベクトルを
列に並べる。$P_0^2 = QD^2Q^{-1}$ と $D^2 = D$ に注目
(例題 \ref{rei:diag-power})。
\item \textbf{問 \ref{q:diag-ex10}}:$A = PDP^{-1}$ とすると
$A^k = PD^kP^{-1} = O$。$D^k$ の対角成分はどうなるか。
\item \textbf{問 \ref{q:diag-ex11}}:転置は積の順序を逆にする
(定理 \ref{thm:transpose})。$({}^tP)^{-1}D\,({}^tP)$ の形に整える。
\item \textbf{問 \ref{q:diag-ex12}}:$A$ を掛けた式から $\lambda_2$ 倍した
式を引くと $\vect{v}_2$ の項が消える。
\end{itemize}
}

% =====================================================
\chapter{対角化の威力 ― 行列のべき乗と漸化式}\label{chap:power}
% =====================================================

\section{$A^n$ の計算}

対角化の最初の報酬は、行列のべき乗が一瞬で計算できることである。
$P^{-1}AP = D$(対角行列)とすると $A = PDP^{-1}$ であり、
\[
A^n = (PDP^{-1})(PDP^{-1})\cdots(PDP^{-1}) = P D^n P^{-1}.
\]
途中の $P^{-1}P$ がすべて消えるからである。そして対角行列のべき乗は
\[
D^n = \mat{\lambda_1 & & \\ & \ddots & \\ & & \lambda_k}^{\,n}
= \mat{\lambda_1^n & & \\ & \ddots & \\ & & \lambda_k^n}
\]
と、対角成分を $n$ 乗するだけである。
$A^{100}$ を $100$ 回掛け算する必要はない。
\textbf{よい座標系に移って計算し、最後に元の座標系に戻す}
――$PDP^{-1}$ の三段構えは、この戦略の式形である。

\section{フィボナッチ数列 ― 本書の伏線を回収する}

\begin{rei}[ビネの公式]
フィボナッチ数列 $F_0 = 0$, $F_1 = 1$, $F_{n+2} = F_{n+1} + F_n$ を考える。
隣り合う $2$ 項を束ねたベクトル $\vect{x}_n = \mat{F_{n+1} \\ F_n}$ を作ると、
漸化式は
\[
\vect{x}_{n+1} = \mat{1 & 1 \\ 1 & 0} \vect{x}_n
\]
という\textbf{行列の反復}になる。$A = \mat{1 & 1 \\ 1 & 0}$ とおけば
$\vect{x}_n = A^n \vect{x}_0$ であり、問題は $A^n$ の計算に帰着した。

固有方程式は $\lambda^2 - \lambda - 1 = 0$、固有値は
\[
\varphi = \frac{1 + \sqrt{5}}{2}, \qquad \psi = \frac{1 - \sqrt{5}}{2}
\]
(黄金比とその相棒)。固有ベクトルは $\vect{v}_\lambda = \mat{\lambda \\ 1}$
と取れる(確かめよ)。相異なる $2$ つの固有値をもつので対角化可能で、
$A^n = PD^nP^{-1}$ を計算して成分を取り出すと(計算は各自):
\[
F_n = \frac{\varphi^n - \psi^n}{\sqrt{5}}.
\]
整数しか生まないはずの漸化式の一般項に $\sqrt{5}$ が現れる、
有名な\textbf{ビネの公式}である。$|\psi| < 1$ より $\psi^n \to 0$ なので、
$F_n$ は $\dfrac{\varphi^n}{\sqrt{5}}$ に最も近い整数であり、
隣接比 $F_{n+1}/F_n$ は黄金比 $\varphi$ に収束する。
\textbf{数列の長期的な振る舞いは、絶対値最大の固有値が支配する}
――この原理は人口モデル、マルコフ連鎖、Google のページランクまで貫いている。
\end{rei}

\begin{itembox}[l]{深掘り:第III部の問 \ref{q:basis2} との合流}
問 \ref{q:basis2} で、漸化式 $a_{n+2} = a_{n+1} + a_n$ を満たす数列全体 $V$ が
$2$ 次元のベクトル空間であることを示した。
このことをビネの公式と突き合わせると、見事な景色が現れる。
等比数列 $(\varphi^n)$ と $(\psi^n)$ は、$\varphi, \psi$ が
$\lambda^2 = \lambda + 1$ の解であることから、どちらも $V$ に属する。
しかも明らかに線形独立。$\dim V = 2$ だから、
\textbf{この $2$ 本は $V$ の基底}であり、$V$ のすべての数列
――フィボナッチ数列も――が
\[
a_n = c_1 \varphi^n + c_2 \psi^n
\]
と一意に書ける。ビネの公式とは「フィボナッチ数列の、
等比数列基底に関する座標表示」だったのである。

さらに深く:数列を一つずらすシフト写像 $S : (a_n) \mapsto (a_{n+1})$ は
$V$ 上の線形変換であり、等比数列はその固有ベクトルである
($S(\lambda^n) = (\lambda^{n+1}) = \lambda \cdot (\lambda^n)$)。
つまり\textbf{等比数列の基底とは、シフト写像を対角化する基底}にほかならない。
高校で「特性方程式」と呼んで天下りに使った道具の正体は、
シフト写像の固有方程式だった。第VI部では、同じ構図が微分方程式で
($e^{\lambda t}$ が微分の固有ベクトルとして)再現される。
\end{itembox}

\section{一般の線形漸化式}

フィボナッチの方法は一般化できる。$k$ 項間漸化式
$a_{n+k} = c_{k-1}a_{n+k-1} + \cdots + c_0 a_n$ は、
連続する $k$ 項を束ねることで $k$ 次正方行列(\textbf{同伴行列})の反復になり、
その固有方程式は特性方程式 $\lambda^k = c_{k-1}\lambda^{k-1} + \cdots + c_0$
に一致する。固有値がすべて異なれば、一般項は等比数列の線形結合で書ける。
(特性方程式が重解をもつ場合は対角化できず、$n\lambda^n$ 型の項が現れる。
その理由はジョルダン標準形で完全に説明される。)

\begin{mondai}\label{q:power}
問 \ref{q:eigen} の $A = \mat{1 & 4 \\ 2 & 3}$ について、
対角化を用いて $A^n$ を求めよ。
\end{mondai}

\begin{mondai}\label{q:rec}
漸化式 $a_{n+2} = a_{n+1} + 2a_n$, $a_0 = 1$, $a_1 = 1$ の一般項を、
等比数列の線形結合として求めよ。
\end{mondai}

% =====================================================
\chapter{ケーリー・ハミルトンの定理と最小多項式}\label{chap:ch}
% =====================================================

\section{ケーリー・ハミルトンの定理}

\begin{teiri}[ケーリー・ハミルトン]
$A$ の固有多項式 $\varphi_A(\lambda)$ に $A$ 自身を代入すると零行列になる:
\[
\varphi_A(A) = O.
\]
たとえば $2$ 次の場合、$\varphi_A(\lambda) = \lambda^2 - (\mathrm{tr}\,A)\lambda + \det A$ なので
\[
A^2 - (\mathrm{tr}\,A)\,A + (\det A)\,E = O.
\]
\end{teiri}

高校数学IIIや旧課程で「ハミルトン・ケーリーの定理」として
$2$ 次の場合に触れた読者もいるだろう。一般の $n$ 次で成り立つ。

この定理の効用は、\textbf{$A$ の高いべきを低いべきで表せる}ことである。
$2$ 次なら $A^2 = (\mathrm{tr}\,A)A - (\det A)E$ を繰り返し使えば、
$A^n$ はつねに $\alpha A + \beta E$ の形に落ちる
(対角化できない行列にも使える点が強み)。
また $\det A \neq 0$ なら、上式を変形して
\[
A^{-1} = \frac{1}{\det A}\bigl((\mathrm{tr}\,A)E - A\bigr)
\]
と、逆行列まで $A$ の多項式で書ける。

\begin{itembox}[l]{深掘り:有名な「一行の誤証明」}
$\varphi_A(\lambda) = \det(\lambda E - A)$ に $\lambda = A$ を代入すれば
$\varphi_A(A) = \det(A - A) = \det O = 0$、はい証明終わり
――という議論は\textbf{誤り}である。どこが誤りか、一度立ち止まって
考えてみてほしい。

種明かし:$\varphi_A(A)$ は「多項式 $\varphi_A(\lambda)$ の係数を保ったまま
$\lambda^k$ を $A^k$ に置き換えた\textbf{行列}」である。
一方 $\det(A - A)$ は\textbf{数}(スカラー)であり、そもそも型が違う。
$\det(\lambda E - A)$ の $\lambda$ は「数を入れる箱」として展開された
多項式であって、行列を入れてよい場所ではない
($\lambda E - A$ の $\lambda E$ に $A$ を「掛ける」のか「代入する」のかが
曖昧になる)。正しい証明は余因子行列を巧みに使う
(あるいは「対角化可能な行列では明らか+極限論法」による)。
\textbf{記号の見かけに引きずられず、式の型(スカラーか、行列か、多項式か)を
確認する}――この習慣は、数学のあらゆる場面で誤りからあなたを守る。
\end{itembox}

\section{最小多項式}

$\varphi_A(A) = O$ により、「$A$ を代入すると消える多項式」は必ず存在する。
その中で最も次数の低いものに名前を付ける。

\begin{teigi}[最小多項式]
$m(A) = O$ を満たす多項式のうち、次数最小で最高次係数 $1$ のものを
$A$ の\textbf{最小多項式}といい、$m_A(\lambda)$ と書く。
\end{teigi}

最小多項式は固有多項式を割り切り、しかも固有多項式と同じ根
(固有値)をすべてもつことが知られている。違いは\textbf{重複度だけ}である。
そして、対角化可能性の最も洗練された判定法が次である。

\begin{teiri}[最小多項式による判定]
$A$ が対角化可能 $\iff$ $m_A(\lambda)$ が重根をもたない。
\end{teiri}

\begin{rei}
$A = \mat{1 & 1 \\ 0 & 1}$ の固有多項式は $(\lambda - 1)^2$。
最小多項式の候補は $\lambda - 1$ と $(\lambda - 1)^2$ だが、
$A - E \neq O$ なので $m_A(\lambda) = (\lambda - 1)^2$。重根をもつから
対角化不可能――第 \ref{chap:diag} 章の結論と一致する。
一方 $E = \mat{1 & 0 \\ 0 & 1}$ も固有多項式は $(\lambda - 1)^2$ だが、
$m_E(\lambda) = \lambda - 1$ で重根なし(そもそも対角行列)。
\textbf{固有多項式が同じでも最小多項式は異なりうる}。最小多項式は
固有多項式より細かく行列を見分ける道具であり、第VI部で
ジョルダン標準形を分類する鍵になる。
\end{rei}

\begin{mondai}\label{q:ch}
$A = \mat{1 & 4 \\ 2 & 3}$(問 \ref{q:eigen} の行列)について、\\
　(1) ケーリー・ハミルトンの定理が成り立つこと($A^2 - 4A - 5E = O$)を
直接計算で確かめよ。\\
　(2) この式を用いて $A^{-1}$ を $A$ と $E$ の線形結合で表せ。
\end{mondai}
